\section{Overview}
\textbf{Note for new users.} If you are unfamiliar with sound analysis, SONAR, and the mathematics behind it I strongly recommend that you read the friendly guides in \hyperref[sec:Appendix_A]{Appendix A} and the relevant chapter in the \href{https://docs.waveops.io/mnw/mnw-guide.pdf}{MNW guide}. You do not need to be a maths wizard to use MACS, but a basic knowledge of the concepts will help you understand what is going on. If you enjoy the detail then \hyperref[sec:Appendix_B]{Appendix B} goes deeper into the concepts. In Mahzel's own words:

\begin{faintquote}
	MACS is a tool provided as an advanced suite for enthusiasts of acoustic analysis. It also provides a custom user-build database where you can slowly build your measurements history, along with all your comments about an encounter. Only had a few seconds of contact? No worries just record it and replay the contact for further analysis for as long as you wish. It includes a database matcher, where you can see if a current contact's markings correspond to something in the database. You may have numerous matches, but in the end, you are in charge of the sonar.
\end{faintquote}

\textbf{Download} MACS for Widows, Linux and MacOS from \href{https://github.com/Mahzel/MACS}{here}.

\subsection{Capabilities}
MACS is an acoustic classification tool designed for use with the naval simulation Modern Naval Warfare (MNW) developed by The Maslas Bros and published by Matrix / Slitherine. MACS does not replace the in game tools.

\textbf{MACS features include:}

\begin{compactitem}
	\item Accepting sound input directly from the PC's sound card or other system capability.
	\item Running alongside MNW without significantly impacting performance.
	\item Fine control of several input and output parameters.
	\item Detailed, fine grain multicolour representations of the analysed sound.
	\item Normalization and noise suppression to manage the impact of background noise.
	\item Allowing recording and replay of sounds for future analysis.
	\item A user-managed database of contacts and measurements (a pre-populated downloadable database is available at the MACS website).
	\item A built in matcher to compare current contacts with those in the database.
\end{compactitem}

\subsection{What MACS does and why would I need it?}

MACS provides six complementary signal-processing views to help you identify and characterise surface or submarine contacts by their acoustic signature. The underlying processing is grounded in real passive sonar digital signal analysis: 

\begin{compactitem}
	\item FFT based spectrograms (LOFARgrams).
	\item DEMON envelope detection (DEMONgrams).
	\item Real cepstrum analysis.
	\item Auto correlation.
	\item Power Spectral Density (PSD) analysis.
	\item ZoomFFT analysis.
\end{compactitem}

Explanations of these and how to use them are described in the following sections. Each processing technique provides a different view of the same signal, when used together they can provide a clearer picture of what is happening around you. The more information you have, the better your chances of detecting and classifying a contact.

\textbf{MACS supplements, it does not replace the in-game tools}
\pagebreak